% Generated from locale/tr/content/lambda-calculus/introduction/computable-lambda.tex; do not hand-edit.


\olfileid[tr]{lam}{int}{lrp}
\olsection{Hesaplanabilir Fonksiyonlar \usetoken{S}{lambda definable}}

\begin{thm}
\ollabel{thm:computable-lambda}
Her hesaplanabilir kısmi fonksiyon !!{lambda definable}.
\end{thm}

\begin{proof}
Her kısmi hesaplanabilir $f$ fonksiyonunun bir $F$ lambda terimi tarafından
!!{lambda defined} olduğunu göstermeliyiz. Kleene normal biçim teoremi
gereği, her ilkel özyinelemeli fonksiyonun bir lambda terimi tarafından
!!{lambda defined} olduğunu ve ardından !!{lambda definable} fonksiyonların
uygun bileşimler ve sınırsız arama altında kapalı olduğunu göstermek
yeterlidir. Her ilkel özyinelemeli fonksiyonun bir lambda terimi tarafından
!!{lambda defined} olduğunu göstermek için başlangıç fonksiyonlarının
!!{lambda definable} olduğunu ve !!{lambda definable} kısmi fonksiyonların
bileşim, ilkel özyineleme ve sınırsız arama altında kapalı olduğunu
göstermek yeterlidir.
\end{proof}

İspatın geri kalanını daha okunaklı kılmak için daha alışılmış bir gösterim
kullanacağız. Örneğin $Mxyz$ yerine $M(x,y,z)$ yazacağız. Bu çağrışımlı olsa
da türsüz lambda kalkülüsündeki terimlere aritelerin iliştirilmediğini
anımsamalısınız; aynı $M$ terimi için $M(x,y)$ yazmak da $M(x,y,z,w)$ yazmak
da aynı ölçüde anlamlıdır. Fakat bu gösterim $M$'yi üç değişkenli bir
fonksiyon olarak ele aldığımızı gösterir ve tanımların arkasındaki amacı
daha açık kılar. Benzer biçimde ``$M$'yi
$M=\lambd[x][\lambd[y][\lambd[z][\dots]]]$ ile tanımlayın'' yerine
``$M$'yi $M(x,y,z)=\dots$ ile tanımlayın'' diyeceğiz.

